\section*{Chapter \ref{ch:b3:gravity-geodesy} --- Gravity and the Figure of the Earth}

\begin{solution}{exo:b3:gravity-geodesy:1}
$b = a(1 - f) = 6378137\times0.99664719 = \qty{6356752.3}{m}$, and
$a - b = \qty{21384.7}{m}$, about \qty{21.4}{km}. Clairaut:
$f = (3J_2 + m)/2 = (0.00324789 + 0.00344979)/2 = \num{0.0033488}$, or
$1/298.61$, against $1/298.26$: \qty{0.12}{\percent} smaller, the size of the
second-order terms the relation neglects.
\end{solution}

\begin{solution}{exo:b3:gravity-geodesy:2}
$\sigma = \qty{6.0e13}{kg}/\qty{5.0e11}{m^2} = \qty{120}{kg.m^{-2}}$, so
$h_w = 120/1000 = \qty{0.12}{m}$: \qty{12}{cm} of water. Delivered to the
ocean, $60/360 = \qty{0.17}{mm}$ of sea level.
\end{solution}

\begin{solution}{exo:b3:gravity-geodesy:3}
At the equator $\gamma = \qty{9.780327}{m.s^{-2}}$. At \ang{45},
$\sin^2\varphi = 0.5$ and $\sin^22\varphi = 1$: $\gamma = 9.780327\times1.0026454
= \qty{9.806200}{m.s^{-2}}$. At the poles, $9.780327\times1.0053024 =
\qty{9.832186}{m.s^{-2}}$. The difference is \qty{0.051859}{m.s^{-2}}, or
\qty{5186}{\milli\gal}.
\end{solution}

\begin{solution}{exo:b3:gravity-geodesy:4}
$D = ET_e^3/(12(1 - \nu^2)) = 10^{11}T_e^3/11.25$: \qty{8.89e21}{N.m} for
\qty{10}{km}, \qty{2.40e23}{N.m} for \qty{30}{km}. With $(\rho_m - \rho_w)g =
\qty{22269}{N.m^{-3}}$, $\alpha = (4D/22269)^{1/4}$: \qty{35.5}{km} and
\qty{81.0}{km}. The moat ends at $3\pi\alpha/4$, \qty{84}{km} and \qty{191}{km};
the \omterm{def:b3:gravity-geodesy:alpha}{forebulge} crests at $\pi\alpha$, \qty{112}{km} and \qty{255}{km}.
\end{solution}

\begin{solution}{exo:b3:gravity-geodesy:5}
(a) $T = -\gamma N = -9.80\times50 = \qty{-490}{J.kg^{-1}}$. (b) The geoid tilts
by $10/100\,000 = \num{1.0e-4}$ rad relative to the ellipsoid, and the plumb
line, normal to the geoid, by the same angle relative to the ellipsoid's normal:
$\num{1.0e-4}\times206\,265 = 20.6''$.
\end{solution}

\begin{solution}{exo:b3:gravity-geodesy:6}
$H = J_2/(C/Ma^2) = \num{1.08263e-3}/0.3308 = \num{0.0032728}$, and
$1/H = 305.55$. $J_2$ is the same by hypothesis; a uniform Earth would have
$C/(Ma^2) = 0.4$ and so $H = J_2/0.4 = \num{0.0027066}$: precession would be
slower. Measuring both $J_2$ and $H$ is what reveals the concentration of mass
towards the centre.
\end{solution}

\begin{solution}{exo:b3:gravity-geodesy:7}
$r = \qty{6878.137}{km}$; $n = \sqrt{GM/r^3} = \qty{1.10678e-3}{rad.s^{-1}}$, a
period of $2\pi/n = \qty{94.6}{min}$. $\dot\Omega = -1.5\times\num{1.10678e-3}
\times\num{1.08263e-3}\times(6378.137/6878.137)^2\times\cos\ang{89} =
\qty{-2.70e-8}{rad.s^{-1}}$, or \ang{-0.134} per day, westward. A near-polar orbit
passes over every latitude, so the whole Earth is mapped; it also precesses
slowly, so that the ground tracks fill in regularly.
\end{solution}

\begin{solution}{exo:b3:gravity-geodesy:8}
$e^{-2\pi\times500/100} = e^{-31.4} = \num{2.3e-14}$: \qty{2e-13}{\milli\gal},
nothing. $e^{-2\pi\times500/1000} = e^{-3.14} = 0.0432$:
\qty{0.43}{\milli\gal}.
\end{solution}

\begin{solution}{exo:b3:gravity-geodesy:9}
$3\pi\alpha/4 = \qty{120}{km}$ gives $\alpha = \qty{50.9}{km}$.
$D = (\rho_m - \rho_w)g\alpha^4/4 = 22269\times(50\,930)^4/4 =
\qty{3.75e22}{N.m}$, and $T_e = (11.25D/E)^{1/3} = \qty{16.2}{km}$. The plate
crosses its unflexed level at $\pi\alpha/2 = \qty{80}{km}$ from the trench.
\end{solution}

\begin{solution}{exo:b3:gravity-geodesy:10}
(a) With $u = x/\alpha$, $\mathrm d^4/\mathrm dx^4 = \alpha^{-4}\,\mathrm
d^4/\mathrm du^4$. Writing $\cos u + \sin u = \sqrt2\cos(u - \pi/4)$, the
function is $\sqrt2w_0\,\mathrm{Re}\,e^{(-1 + i)u - i\pi/4}$, and
$(-1 + i)^4 = -4$: so $w'''' = -4w/\alpha^4$. Then $Dw'''' + (\rho_m -
\rho_i)gw = (-4D/\alpha^4 + (\rho_m - \rho_i)g)w = 0$, since $\alpha^4 =
4D/((\rho_m - \rho_i)g)$. (b) $w'' = -(2w_0/\alpha^2)e^{-u}(\cos u - \sin u)$, so
$M = -Dw'' = (2Dw_0/\alpha^2)e^{-u}(\cos u - \sin u)$: largest in magnitude at
$u = 0$, zero where $\tan u = 1$, $x = \pi\alpha/4$. (c) Under the load the
plate is concave upward, its upper surface in compression; beyond $\pi\alpha/4$
it is convex upward, its upper surface stretched. There $M' \propto
-e^{-u}\cos u$ vanishes at $u = \pi/2$: the stretching is greatest at
$x = \pi\alpha/2$, on the outer slope of the moat towards the \omterm{def:b3:gravity-geodesy:alpha}{forebulge}, where
$|M|$ is $e^{-\pi/2} = \qty{21}{\percent}$ of its value under the load. That is
where normal faults and cracks would open at the top of the plate.
\end{solution}

\begin{solution}{exo:b3:gravity-geodesy:11}
$\sqrt{1 + \eta} = 2.5(1 - 1.5\times0.3308) = 1.2595$, so $1 + \eta = 1.5863$
and $\tfrac52m/f = 2.5863$: $f_h = 2.5\times0.00344979/2.5863 = \num{0.0033346}$,
$1/f_h = 299.9$. The difference $f - f_h = \num{1.82e-5}$; times $a$, the
equator stands \qty{116}{m} above the hydrostatic figure. The Earth is not quite
a fluid in equilibrium: its interior sustains stresses, from the strength of
the lithosphere or the flow and density differences of the mantle, that hold
the shape a little away from it.
\end{solution}

\begin{solution}{exo:b3:gravity-geodesy:12}
The node must advance by $2\pi/(365.242\times86\,400) =
\qty{1.991e-7}{rad.s^{-1}}$. At $r = \qty{7078.137}{km}$, $n =
\qty{1.0602e-3}{rad.s^{-1}}$ and $\tfrac32nJ_2(a/r)^2 = \num{1.398e-6}$, so
$\cos i = -\num{1.991e-7}/\num{1.398e-6} = -0.1424$ and $i = \ang{98.2}$: a
slightly retrograde, sun-synchronous orbit.
\end{solution}

\begin{solution}{pb:b3:gravity-geodesy:1}
\textbf{1.} $m = (\num{7.292115e-5})^2\times6378137^2\times6356752.3/
\num{3.986005e14} = \num{0.00344979}$.

\textbf{2.} $f = (3J_2 + m)/2 = \num{0.0033488}$, $1/f = 298.61$.

\textbf{3.} $\sqrt{1 + \eta} = 2.5(1 - 1.5\times0.3308) = 1.2595$,
$\tfrac52m/f_h = 2 + 0.5863$, so $f_h = \num{0.0033346}$ and $1/f_h = 299.88$.

\textbf{4.} The Clairaut value, $1/298.61$, is within \qty{0.12}{\percent} of
the observed $1/298.26$; the hydrostatic one, $1/299.88$, is smaller by
$\Delta f = \num{1.82e-5}$, and $a\,\Delta f = \qty{116}{m}$: the equator stands
\qty{116}{m} above the hydrostatic figure.

\textbf{5.} $H = \num{1.08263e-3}/0.3308 = \num{0.0032728}$ ($1/305.6$).

\textbf{6.} $\tfrac12\times\qty{80}{km}\times\qty{6}{km} = \qty{240}{km^2} =
\qty{2.40e8}{m^2}$.

\textbf{7.} $V_0 = (2800 - 1030)\times9.81\times\num{2.40e8} =
\qty{4.17e12}{N.m^{-1}}$.

\textbf{8.} Airy: $(\rho_m - \rho_v)r = (\rho_v - \rho_w)h$, so $r =
6\times1770/500 = \qty{21.2}{km}$ under the crest.

\textbf{9.} The sea floor is deflected far beyond the foot of the ridge: the
load is not compensated column by column but by the bending of the plate over a
region four times wider. The plate has strength: regional compensation.

\textbf{10.} $\alpha = 165/(3\pi/4) = \qty{70.0}{km}$.

\textbf{11.} $\alpha = 220/\pi = \qty{70.0}{km}$: the same, as a single elastic
plate requires.

\textbf{12.} $D = (\rho_m - \rho_w)g\alpha^4/4 = 2270\times9.81\times(\num{7.00e4})^4/4
= \qty{1.34e23}{N.m}$.

\textbf{13.} $T_e = (12(1 - \nu^2)D/E)^{1/3} = (11.25\times\num{1.34e23}/
\num{1e11})^{1/3} = \qty{24.7}{km}$.

\textbf{14.} $w_0 = V_0\alpha^3/(8D) = \num{4.17e12}\times(\num{7.00e4})^3/
(8\times\num{1.34e23}) = \qty{1334}{m}$.

\textbf{15.} $w_0e^{-\pi} = 1334\times0.0432 = \qty{58}{m}$.

\textbf{16.} $2\times1334\times\num{7.00e4}\times2270\times9.81 =
\qty{4.16e12}{N.m^{-1}}$, equal to $V_0$ to within the rounding of $D$.

\textbf{17.} \qty{24.7}{km} lies inside the range of the Hawaiian Ridge. The
plate sinks \qty{1.3}{km} under the axis instead of growing a root of
\qty{21}{km}: the same buoyancy is spread over the \qty{330}{km} of the moat.

\textbf{18.} $\lambda = 2\pi\times\qty{70.0}{km} = \qty{440}{km}$,
$k = 1/\alpha = \qty{1.43e-5}{m^{-1}}$.

\textbf{19.} $\sigma = (2900 - 1030)\times500 = \qty{9.35e5}{kg.m^{-2}}$, so
$2\pi G\sigma = \qty{3.92e-4}{m.s^{-2}} = \qty{39.2}{\milli\gal}$ at the sea floor;
times $e^{-k\times\qty{5}{km}} = 0.931$, \qty{36.5}{\milli\gal} at sea level
(negative where the plate is pushed down).

\textbf{20.} $\sigma = (3300 - 2900)\times500 = \qty{2.0e5}{kg.m^{-2}}$:
\qty{8.39}{\milli\gal}, times $e^{-k\times\qty{12}{km}} = 0.843$,
\qty{7.07}{\milli\gal}. The total is \qty{43.6}{\milli\gal}.

\textbf{21.} $e^{-k\times\qty{400}{km}} = e^{-5.71} = 0.0033$: \qty{0.14}{\milli\gal}.

\textbf{22.} $e^{-2\pi\times400/80} = e^{-31.4} = \num{2.3e-14}$: the satellite
sees nothing of the ridge itself, only the broad flexure, and even that
weakened three-hundred-fold; the ridge's own anomaly must be measured at sea.

\textbf{23.} $n \approx 2\pi\times6371/440 = 91$.

\textbf{24.} From questions 10 to 13, the \omterm{def:b3:gravity-geodesy:flexure}{effective elastic thickness} of the
plate beneath the Oru Ridge is
\textbf{\qty[text-series-to-math=true]{24.7}{km}}.
\end{solution}
